\usepackage[most]{tcolorbox}
\tcbset{
  defstyle/.style={
    enhanced,
    colback=blue!5!white,
    colframe=blue!75!black,
    fonttitle=\bfseries,
    title=Definition,
    sharp corners,
    boxrule=0.8pt,
    breakable,
  }
}

\newenvironment{definitionbox}[1][]{
  \begin{tcolorbox}[defstyle,title={#1}]
}{
  \end{tcolorbox}
}

% Usage
\begin{definitionbox}[14.3 Definition — Cauchy Criterion]
We say a series $\sum a_n$ satisfies the \emph{Cauchy criterion} if its
sequence $(s_n)$ of partial sums is a Cauchy sequence:
\[
\text{for each } \varepsilon>0, \text{ there exists an } N
\text{ such that } m,n > N \implies |s_n - s_m| < \varepsilon.
\]
\end{definitionbox}


MORE POLISHED DESIGN:

\tcbset{
  fancystyle/.style={
    enhanced,
    colback=blue!3!white,
    colframe=blue!60!black,
    fonttitle=\bfseries,
    title=Definition,
    attach boxed title to top left={yshift=-2mm,xshift=2mm},
    boxed title style={colback=blue!60!black, colframe=blue!60!black},
    coltitle=white,
    rounded corners,
    shadow={0.7mm}{-0.7mm}{0mm}{black!20!white},
    breakable,
  }
}

\newenvironment{fancydef}[1][]{
  \begin{tcolorbox}[fancystyle,title={#1}]
}{
  \end{tcolorbox}
}

% Usage
\begin{fancydef}[14.3 Definition — Cauchy Criterion]
We say a series $\sum a_n$ satisfies the \emph{Cauchy criterion} if its
sequence $(s_n)$ of partial sums is a Cauchy sequence.
\end{fancydef}